\documentclass{subfiles}
\begin{document}
    As we shall see, insertion and deletion on a RB tree are everything but simple. 
        We shall in fact see that either of the two operations
        may cause one or more properties to fail.
        To fix this issues specific procedures are used to restore each property;
        in doing so some local pointer manipulations are made,
        which here we refer to as rotations.
        We describe left rotations, the other being symmetrical.

    \subfile{../TikZ/Figure 1 - RB rotations}
    Consider the case shown in \Cref{fig:RB-rotations-a}.
        To obtain the tree in \Cref{fig:RB-rotations-b} we proceed as follow:
        we let \(y\)'s left child to be \(x\)'s new right child,
        and we make \(x\) be the new \(y\)'s left child.
\end{document}
